\subsection{The full group in equal characteristic: completion of B3}

In this subsection $R$ is a henselian DVR of equal characteristic $p>0$,
$K=\operatorname{Frac}(R)$, and $k$ is its residue field. The conditions
$a,b,c,d,e$ have the meanings in B3, with the \emph{full} geometric
fundamental group in $b,c,e$, and $\ell\ne p$ in $d$. A residue splitting
is assumed given whenever $e$ is mentioned. These are modern extensions
of cote 149, manuscript pp.~49--51, scans 99--104; we do not attribute
the positive-characteristic conclusions below to the manuscript.

\begin{claim}{B.20}{Proved}
For a proper hyperbolic curve $C/K$, trivial full outer inertia implies
that $C$ has good reduction and that the full geometric specialization map
\[
 \operatorname{sp}:\pi_1(C_{\bar K})\longrightarrow
                  \pi_1(\mathcal C_{\bar k})
\]
is an isomorphism. Its outer equivariance identifies the action restricted
along \emph{every} residue splitting with that of the smooth special
fibre. In particular $b\Rightarrow e$.
For an affine hyperbolic curve in equal characteristic $p$, full inertia
acts nontrivially already on $H^1_{\mathrm{et}}(U_{\bar K},\F_p)$.
Consequently full $b$ and $c$ are impossible in the affine case.
\end{claim}
\begin{proof}
We first prove the proper assertion. Trivial full outer inertia implies
triviality on every characteristic pro-$\ell$ quotient. Tamagawa's
good-reduction criterion therefore gives a smooth proper model
$\mathcal C/R$ \citep[Theorem 5.3]{BC_Tamagawa}. The specialization map
is surjective: SGA 1, X, Corollary 2.4 applies to every proper smooth
morphism over a locally noetherian base, not just to a complete trait
\citep[X, Corollary 2.4]{BC_SGA1}. We prove its injectivity by treating
\emph{all} finite geometric Galois covers, including those of degree
divisible by $p$.

Let $D_{\bar K}\to C_{\bar K}$ be a connected finite etale Galois cover
with group $Q$ and degree $d$. Finite presentation allows us, after a
finite separable extension $L/K$, to descend the cover, all its deck
automorphisms and a point of $D$ to $L$. The latter gives compatible
sections of the arithmetic sequences of $D$ and $C$, and hence actual
Galois actions on $H=\pi_1(D_{\bar K})\triangleleft
\Pi=\pi_1(C_{\bar K})$. For $\sigma\in I_L$, the action on $\Pi$ is
inner by hypothesis. The induced outer action on $H$ consequently belongs
to the finite image of $\Pi/H$ in $\Out(H)$: conjugation by an element
of $H$ is inner on $H$. This conclusion does not require a choice of
continuous conjugating elements in $\Pi$.

The continuous outer action of $I_L$ on $H$ thus has finite image. Its
kernel is open in $I_L$; choose an open subgroup of $G_L$ whose
intersection with $I_L$ lies in that kernel. After replacing $L$ by the
corresponding finite extension, the outer inertia action on $H$, and
therefore on its pro-$\ell$ quotient, is trivial. Applying the same
good-reduction criterion to the proper curve $D$ gives a smooth proper
model $\mathcal D/R_L$. Here $g(D)=1+d(g(C)-1)\ge2$.

We now justify extension of the \emph{map}; good models alone would not
suffice. The chosen $L$-point of $D$ extends by properness and gives an
$R_L$-section of $\mathcal D$ and, by its image on the generic fibre,
an $R_L$-section of $\mathcal C$. The relative Abel map embeds
$\mathcal C$ as a closed subscheme of its relative Jacobian $J$.
The Neron mapping property extends the generic map
$\mathcal D_L\to\mathcal C_L\hookrightarrow J_L$ uniquely across
the smooth $R_L$-scheme $\mathcal D$.
For precise sources, SGA 7 I, IX, 1.1 defines that mapping property on
\emph{all} smooth schemes over the trait, and Theorem 12.1(c)--(d),
especially (12.1.8)--(12.1.10), identifies the relative degree-zero Picard
scheme with the identity component of the Neron model
\citep[IX, 1.1 and Theorem 12.1]{BC_SGA7}. One may apply the mapping
property to the Neron model first: the extended section maps to zero,
and the connected special fibre of $\mathcal D$ then maps into its
identity component $J$. These identifications can be checked over the
strict henselization and descend by their uniqueness.
The extended map factors through the closed subscheme $\mathcal C$:
the pulled-back defining ideal vanishes on the dense generic fibre and
$\mathcal O_{\mathcal D}$ is torsion-free over the DVR. We obtain
$f:\mathcal D\to\mathcal C$.

Choose a relatively ample line bundle $M$ on $\mathcal C$. The degree
of $f^*M$ on fibres of the proper flat family $\mathcal D/R_L$ is
constant, by constancy of its Euler characteristic. Its generic degree
is positive. Therefore the map on the special fibre is nonconstant.
All its fibres are finite, so $f$ is proper and quasi-finite, hence finite.
At a point of $\mathcal D$ above a point of $\mathcal C$, both local
rings are regular and have the same dimension. A system of parameters
from the latter is a system of parameters in the former because the
map is finite. Since a regular local ring is Cohen--Macaulay, this is
a regular sequence, and the local flatness criterion gives flatness.
Thus $f$ is finite flat of rank $d$.

Every generic deck automorphism extends uniquely to $\mathcal D$.
Here the exact input is Deligne--Mumford's Theorem 1.11: the isomorphism
scheme of two stable curves is finite and unramified; its properness
extends the automorphism over the DVR
\citep[Theorem 1.11, p.~84]{BC_DM}. Specialization of these automorphisms
is injective. Indeed the equalizer of two sections of an unramified
separated scheme is open and closed; if it contains the closed point
of a local trait, it is the entire trait. They commute with $f$ by
uniqueness on the dense generic fibre. The special map, of total degree
$d=|Q|$, consequently has $d$ distinct deck automorphisms. Its separable
degree is therefore at least $d$, so it is a separable Galois map of
degree $d$ and has no inseparable factor.

The genera of the smooth proper fibres are constant, and the generic
cover is etale. Hence
\[
 2g(\mathcal D_k)-2=d\bigl(2g(\mathcal C_k)-2\bigr).
\]
Riemann--Hurwitz on the special fibres makes its effective different
zero. Thus the special map is etale. A finite flat morphism which is
etale on both fibres of a trait is etale; this follows either from
the fibrewise criterion or from the vanishing of its module of
relative differentials. We have proved that the given geometric cover
extends etale after a finite extension of the trait.

Specialization is compatible with that extension and with this finite
etale cover. Accordingly the quotient $\Pi\to Q$ factors through
$\pi_1(\mathcal C_{\bar k})$. The change of algebraically closed fields
does not change these \emph{proper} geometric groups
\citep[X, Corollary 1.8]{BC_SGA1}. Every finite quotient of $\Pi$ has
now been treated; their kernels have trivial intersection, proving
injectivity of $\operatorname{sp}$. The construction of specialization
from the arithmetic fundamental-group diagrams in SGA 1, X, 2.1--2.4
also proves its outer $G_K$-equivariance, with $G_K$ acting on the
special group through $G_k$. Since $\operatorname{sp}$ is an isomorphism,
restriction along any $s:G_k\to G_K$ is exactly the special-fibre
outer action. This proves $e$ over $k$ itself.

For the affine assertion, passing to a finite separable extension
preserves any putative triviality of inertia. We may therefore split a
boundary point $P$ and choose, by Riemann--Roch, a function
$f\in\mathcal O(U)$ whose only pole is at $P$, of order $m>0$ prime
to $p$. We choose $m$ sufficiently large. If $\varpi$ is a uniformizer,
let $a$ satisfy $a^p-a=\varpi^{-1}$. This defines a totally ramified
Artin--Schreier extension of degree $p$: the valuation of $a$ is
$-1/p$, so the extension has ramification index $p$. Some
$\tau\in I_K$ satisfies $\tau(a)=a+1$. On the affine geometric curve,
the Artin--Schreier sequence gives
\[
 H^1_{\mathrm{et}}(U_{\bar K},\F_p)
 =\mathcal O(U_{\bar K})/(h^p-h).
\]
Then $\tau[af]-[af]=[f]\ne0$. To check the inequality, if
$h^p-h=f$, any pole of $h$ gives a pole of $f$ of order divisible by
$p$, and no pole can disappear. The unique pole of $f$ has order $m$
prime to $p$, a contradiction. Thus the inertia action on the mod-$p$
abelian quotient is nontrivial, proving the last assertion.
\end{proof}

\begin{claim}{B.21}{Proved}
Let $k=\F_9=\F_3(\alpha)$ with $\alpha^2=-1$, $R=k[[t]]$,
$K=k((t))$, and $\lambda=-1+t$. The smooth proper genus-two curve
\[
 C:\quad y^2=(x^2+\alpha)(x^2+\alpha-1)(x^2+\alpha-\lambda)
\]
has $a,d$ true and $b,c,e$ false. Failure of $e$ holds for every
residue splitting and after every finite residue extension.
\end{claim}
\begin{proof}
Consider the smooth elliptic model
\[
 E:\quad v^2=u(u-1)(u-\lambda)=u^3-tu^2+\lambda u.
\]
Its discriminant $16\lambda^2(1-\lambda)^2$ is a unit. For a
characteristic-three cubic $v^2=u^3+a_2u^2+a_4u+a_6$, the nonzero
three-torsion points are the flexes, with
$a_2u^3+a_2a_6-a_4^2=0$. This is also the explicit calculation of
\citet[Example 5.27]{BC_MoonenIsogs}. In our case their $u$ coordinates
satisfy $u^3=-\lambda^2/t$.

We use this formula on the \emph{etale Verschiebung kernel}, avoiding
the imperfection pitfall in substituting $K^{\mathrm{sep}}$ for $\bar K$
in $E[3]$. Let
$F:E\to E^{(3)}$ be relative Frobenius and
$V:E^{(3)}\to E$ Verschiebung, so $VF=[3]$; these maps and the
identity over a characteristic-$p$ base are constructed in
\citet[Definitions 5.18 and Propositions 5.19--5.20]{BC_MoonenIsogs}.
The two nonzero points of $\ker V$ have coordinates
\[
 X=-\frac{\lambda^2}{t},\qquad
 Y^2=X(X-1)(X-\lambda^3)
     =-\frac{\lambda^4(\lambda-1)^4}{t^3}.
\]
Indeed Frobenius is radicial and bijective on algebraically closed
points, and carries the two flexes to these points; $\deg V=3$.
On putting $z=Yt^2/(\lambda^2(\lambda-1)^2)$ we obtain
\[
 z^2=-t.
\]
This quadratic polynomial is separable and Eisenstein. Thus the two
nonzero points of the \emph{finite etale} $K$-group scheme $\ker V$
are exchanged by inertia. Notice that the assertion concerns a cover
$E^{(3)}\to E$ which is etale over $K$, not a possibly nonetale
$3$-torsion subgroup of $E$ over this imperfect field; the distinction
is explicitly warned about in \citet[Caution 5.24]{BC_MoonenIsogs}.

The map $C\to E$, $(x,y)\mapsto(x^2+\alpha,y)$, has degree two.
The displayed sextic expands over $\F_9[[t]]$ as
\[
 f(x)=x^6-tx^4+\bigl((\alpha+1)t-1\bigr)x^2
                  +\alpha+(\alpha+1)t.
\]
At $t=0$ its derivative is $x$, while $f(0)=\alpha\ne0$.
The two points at infinity are smooth since the leading coefficient
is one and $2$ is invertible. Hence the proper model is smooth of genus
two, proving $a$ and $d$.
Pulling back the degree-three etale Verschiebung cover along $C\to E$
gives a geometrically connected cover: the two function-field degrees
two and three are coprime, and the latter extension is Galois over
the algebraic closure. Thus the geometric abelianized group of $C$
has the inertia-ramified quotient $\ker V(\bar K)$.
It follows that $c$, and therefore $b$, is false.

We compute the geometric $3$-ranks to exclude $e$. On the holomorphic
basis $\eta_0=dx/y$, $\eta_1=x\,dx/y$, Cartier's formula gives
\[
 \mathcal C(\eta_0)=\bigl((\alpha+1)t-1\bigr)^{1/3}\eta_0,
 \qquad
 \mathcal C(\eta_1)=(-t)^{1/3}\eta_1.
\]
Here we compute over an algebraic closure: write $1/y=f(x)/y^3$
and use $\mathcal C(x^{3j+2}dx)=x^jdx$. The semilinear map is invertible
generically and has stable rank one at $t=0$. Thus
$f(C_{\bar K})=2$ and $f(C_{\bar k})=1$.

Suppose the full residue action were realized by a smooth hyperbolic
curve $W$ over a finite extension $k'/k$. The intrinsic pro-prime-to-$3$
quotient then has the same outer action as $C_{k'}$: good proper
specialization identifies it independently of the splitting. Cohomology
at any $\ell\ne3$ distinguishes a proper curve from an affine curve
and determines its genus; hence $W$ is proper of genus two. Centre-free
extension reconstruction B.2 identifies the associated geometrically
prime-to-$3$ arithmetic groups. Saidi--Tamagawa's exact finite-field
anabelian theorem applies: their Theorem 3.9 assumes all primes other
than the field characteristic in the geometric group and concludes that
an isomorphism of those arithmetic groups comes from a scheme
isomorphism \citep[Theorem 3.9]{A_SaidiTamagawa2009}.
Consequently $W$ is isomorphic to $C_{k'}$ as a scheme; a possible
automorphism of its finite field of constants does not affect its
$3$-rank. But its full geometric group was assumed isomorphic to
$\pi_1(C_{\bar K})$, and
$\dim_{\F_3}H^1(\pi_1(-),\F_3)$ is the $3$-rank. It would have
to be both two and one. This proves failure of $e$, for all the
splittings and finite residue extensions asserted.
\end{proof}

\begin{claim}{B.22}{Proved}
Over $\F_7((t))$ the proper genus-two curve
\[
 D:\quad y^2=(x^3-t^6)(x^3-1)
\]
has $c,d$ true and $a,b,e$ false for the full-group conditions.
\end{claim}
\begin{proof}
The normalized compact-type model is the same explicit model as in
B.11, with uniformizer $t$. We recall its checks. In the inner chart
$x=t^2X$, $y=t^3Y$, the special component is $Y^2=1-X^3$.
The outer component normalizes to $V^2=x(x^3-1)$, which becomes
$Y^2=1-X^3$ under $X=1/x$, $Y=V/x^2$. Near the connecting point put
$xz=t^2$ and $A=(1-z^3)(x^3-1)$. The normalization has
\[
 V^2=xA,\qquad W^2=z/A,\qquad VW=t.
\]
Since $A(0,0)=-1$, this is a regular separating node of thickness one.
Both components are smooth elliptic curves in characteristic seven,
and the node gives a stable compact-type special fibre. The model
therefore remains nodal after every finite extension, so $a$ is false
and Tamagawa makes $b$ false. C.5 excludes even prime-to-$7$ residue
realization, hence excludes full $e$.

The relative Jacobian $J$ is an abelian scheme with special fibre the
product of the two component Jacobians. Each is ordinary: the Hasse
coefficient of $x^6$ in $(1-x^3)^3$ is $3\ne0$ in $\F_7$.
For $\ell\ne7$ good reduction gives trivial inertia on $T_\ell J$.
For the residue prime we give the additional argument rather than
appealing to a prime-to-$p$ criterion. The iterated Verschiebung
\[
 V^n:J^{(7^n)}\longrightarrow J
\]
is finite flat, of rank $7^{2n}$, and its special kernel is etale because
the special abelian surface is ordinary. Thus its kernel is finite
etale over $R$: the nonetale locus of this finite flat scheme is closed,
its image in the local trait is closed, and it misses the closed point.
Over an algebraically closed field, relative Frobenius is a bijection
on points and identifies these kernels with the etale $7^n$-torsion
quotients. They consequently give unramified $7$-primary geometric
Tate modules. More directly, the etale isogenies $V^n$ form the
cofinal $7$-power covers of the ordinary abelian variety, all extended
over $R$. Hence inertia is trivial on every primary part of
$\pi_1(D_{\bar K})^{\mathrm{ab}}=\pi_1(J_{\bar K})$.
This proves $c$ and $d$.
\end{proof}

\begin{claim}{B.23}{Proved}
Over $K=\F_3((t))$ the proper genus-two curve
\[
 C:\quad y^2=t(x^5-x)
\]
has $e$ true for every residue splitting, whereas $a,b,c,d$ are false.
\end{claim}
\begin{proof}
The constant curve $C_0:Y^2=x^5-x$ over $\F_3$ is smooth of genus two:
the polynomial and its derivative are coprime. Over $L=K(r)$,
$r^2=t$, the substitution $y=rY$ identifies $C_L$ with $(C_0)_L$.
The nontrivial inertia element acts on $C_0$ as the hyperelliptic
involution $\iota$. It acts as $-1$ on every prime-to-$3$ Tate module
of the Jacobian: the quotient is $\mathbf P^1$, so
$1+\iota=0$ on its degree-zero Picard variety. Thus $d,c,b$ fail,
and good reduction would contradict $a\Rightarrow d$.

Because the curve is \emph{proper}, base extension from an algebraic
closure of the finite field to $\bar K$ preserves its full geometric
fundamental group. Its full outer $G_K$-action is the constant
$G_{\F_3}$-action twisted by the quadratic character
$\chi:G_K\to\langle\iota\rangle$. For any given splitting $s$,
the character $\chi\circ s$ defines a quadratic twist of $C_0$ over
$\F_3$; ordinary finite Galois descent makes this a smooth proper
genus-two curve. Its full outer action is exactly the restriction along
$s$. This proves $e$ already over $\F_3$. Properness in this argument
is essential; we do not substitute an affine constant curve and assume
invariance of its full wild fundamental group under extension of
algebraically closed fields.
\end{proof}

\begin{claim}{B.24}{Proved}
The same failure of $a\Rightarrow b,c,e$ occurs explicitly at $p=2$.
Over $R=\F_2[[t]]$, let $C$ be the smooth proper normalization of
\[
                    y^2+t z^3y+y=z^9.
\]
Then $C_K$ is a proper genus-four curve with $a,d$ true and full
$b,c,e$ false. Failure of $e$ again holds for every residue splitting
and after every finite residue extension. The ramified character which
detects $c$ here occurs on a cyclic quotient of order four.
\end{claim}
\begin{proof}
Start with the smooth elliptic scheme
\[
 E:\quad y^2+txy+y=x^3.
\]
Its discriminant is $1+t^3$, a unit. Its generic fibre is ordinary
since its Weierstrass coefficient $a_1=t$ is nonzero, whereas its special
fibre $y^2+y=x^3$ is supersingular; the elementary characteristic-two
criterion is \citet[Example 5.26]{BC_MoonenIsogs}.
The divisor of $x$ is $P+Q-2O$, where $P=(0,0)$, $Q=(0,1)$,
and $O$ is the origin at infinity. These are three disjoint sections.
Adjoining $z^3=x$ gives the displayed curve. The normalization is smooth
over $R$: away from the three sections the cover is etale, and at each
section its tame cyclic local normalization, after an etale change
extracting a cube root of the unit, is a smooth parameter with base
parameter its cube. All three ramification indices are three, and
Riemann--Hurwitz gives $2g(C)-2=3(3-1)=6$.

We compute a ramified etale cover of the generic elliptic curve. Put
\[
 x=t^2X+t^{-1},\qquad y=t^3Y+t^{-3},\qquad
 B=t^{-3},\qquad A=t^{-12}+t^{-9}.
\]
Its generic equation becomes
\[
                 Y^2+XY=X^3+BX^2+A.                         \tag{B.24.1}
\]
For a point $(X,Y)$ with $X\ne0$, the tangent slope is
$X+Y/X$. Substitution in the duplication formula gives
\[
                  X(2P)=X(P)^2+A/X(P)^2.
\]
The unique nonzero geometric two-torsion point has $X=0$.
Consequently the points of exact order four have $X^4=A$.
Under twice iterated relative Frobenius their images are the two
primitive points of the etale kernel of
$V^2:E^{(4)}\to E$. Writing their coordinates as $(A,Y_4)$, we have
\[
 Y_4^2+AY_4=A^3+B^4A^2+A^4.
\]
Set $w=Y_4/A+A+B^2+B$. The equation is exactly
\[
                       w^2+w=B=t^{-3}.                       \tag{B.24.2}
\]
The right side has odd pole order three, so this Artin--Schreier
extension is totally ramified of degree two: any root has valuation
$-3/2$. It exchanges the two primitive generators of the etale cyclic
group $\ker V^2(\bar K)\simeq\Z/4$. The relative Frobenius map is
used only to locate this etale kernel; no identification of all
geometric torsion points with points over $K^{\mathrm{sep}}$ is made.
Its pullback along the degree-three map $C\to E$ is geometrically
connected, since three and four are coprime. Thus it gives an
inertia-ramified quotient of $\pi_1(C_{\bar K})^{\mathrm{ab}}$.
This proves failure of $c$ and $b$, despite good reduction.

For completeness, the special curve $y^2+y=z^9$ has $2$-rank zero
directly by Cartier. The different exponent at its unique point over
$z=\infty$ is ten, so $dz$ has divisor $6P_\infty$; a basis of
regular differentials is $dz,z\,dz,z^2dz,z^3dz$. Cartier kills the
even powers and sends $z\,dz$ to $dz$ and $z^3dz$ to $z\,dz$.
Its third iterate is zero. Generically, pullback and norm for the
degree-three map to the ordinary $E$ have composite multiplication by
three, so the etale two-primary Tate module of $E$ is a direct summand
of that of $\Jac(C)$. In particular $f(C_{\bar K})\ge1$.
The special $2$-rank is zero, whereas the generic one is positive.
If a full residue realization existed, prime-to-two specialization,
the properness/genus invariants and Saidi--Tamagawa Theorem 3.9 would
identify it with the special curve, exactly as in B.21. Full mod-two
homology contradicts the two different ranks. This proves the claim.
\end{proof}

The exact identities, Hasse coefficients, smoothness and node checks in
B.21--B.24 are verified in \texttt{code/BC\_equal\_char\_checks.py};
all 35 named checks pass. They verify the explicit examples, whereas
the all-cover specialization argument is proved in B.20.
